% Note de rendu - test Buddimm « qualification des permis de construire (Sitadel) »
% Compilation : xelatex note_buddimm.tex (deux passes)
\documentclass[10pt,a4paper]{article}

% ---------- polices et langue ----------
\usepackage{fontspec}
\setmainfont{Libertinus Serif}
\setsansfont{Libertinus Sans}
\usepackage[french]{babel}
\usepackage{microtype}

% ---------- couleurs ----------
\usepackage{xcolor}
\definecolor{encre}{HTML}{1B2438}
\definecolor{accent}{HTML}{2563B0}
\definecolor{vert}{HTML}{1B8059}
\definecolor{orange}{HTML}{C9731F}
\definecolor{rouge}{HTML}{B4443C}
\definecolor{gris}{HTML}{5F6670}
\definecolor{filet}{HTML}{D5DAE1}
\definecolor{fond}{HTML}{F3F5F8}

% ---------- mise en page ----------
\usepackage[a4paper,top=13mm,bottom=16mm,left=17mm,right=17mm,footskip=8mm]{geometry}
\usepackage[hidelinks]{hyperref}
\usepackage{titlesec}
\titleformat{\section}{\sffamily\bfseries\normalsize\color{encre}}
  {\colorbox{accent}{\makebox[3.2mm]{\color{white}\footnotesize\thesection}}}{6pt}{\MakeUppercase}
  [\vspace{-4pt}{\color{filet}\rule{\linewidth}{.5pt}}]
\titlespacing*{\section}{0pt}{9pt}{-1pt}
\usepackage{enumitem}
\setlist{nosep,leftmargin=1.05em,topsep=1pt,itemsep=2.2pt}
\setlist[itemize]{label={\color{accent}\scriptsize$\blacksquare$}}
\usepackage{amssymb}
\usepackage{tabularx,booktabs,array}
\renewcommand{\arraystretch}{1.15}
\newcolumntype{Y}{>{\raggedright\arraybackslash}X}
\newcolumntype{P}[1]{>{\raggedright\arraybackslash}p{#1}}
\setlength{\parindent}{0pt}
\setlength{\parskip}{3pt}
\usepackage{fancyhdr}
\pagestyle{fancy}\fancyhf{}\renewcommand{\headrulewidth}{0pt}
\fancyfoot[L]{\sffamily\scriptsize\color{gris}\textbf{BAQUA Abdellah} · Buddimm, stage IA agentique 2027 · Qualification des demandeurs de permis (Sitadel, Yvelines)}
\fancyfoot[R]{\sffamily\scriptsize\color{gris}\href{https://github.com/AbdellahBaqua/SitadelAgent}{github.com/AbdellahBaqua/SitadelAgent} · \thepage/2}

% ---------- schémas ----------
\usepackage{tikz}
\usetikzlibrary{positioning,arrows.meta,calc,fit,backgrounds}
\tikzset{
  etape/.style={draw=#1,fill=#1!6,rounded corners=2pt,line width=.6pt,align=center,
                font=\sffamily\scriptsize,inner sep=3pt,minimum height=11mm,text width=27mm,text=encre},
  fleche/.style={-{Stealth[length=2.2mm]},line width=.7pt,color=gris},
  chiffre/.style={font=\sffamily\bfseries\scriptsize,color=#1,fill=white,inner sep=1pt},
}
\newcommand{\nb}[1]{\textbf{#1}}
% étiquette + valeur du bloc d'identité
\newcommand{\champ}[2]{{\sffamily\tiny\bfseries\color{gris}\addfontfeature{LetterSpace=8}#1}\\[1.5pt]#2}
% chiffre clé
\newcommand{\kpi}[3]{\begin{minipage}[t]{.235\linewidth}\sffamily
  {\color{#1}\rule{\linewidth}{1.6pt}}\\[2pt]
  {\color{#1}\LARGE\bfseries #2}\\[1pt]
  {\scriptsize\color{encre}#3\par}\end{minipage}}
% colonne établi / fragile / manque
\newcommand{\bilan}[3]{\begin{minipage}[t]{.315\linewidth}
  {\color{#1}\rule{\linewidth}{1.6pt}}\\[1pt]
  {\sffamily\bfseries\small\color{#1}#2}\\[1pt]
  {\footnotesize\raggedright #3\par}\end{minipage}}

\begin{document}

% ================= en-tête =================
{\sffamily
{\color{accent}\scriptsize\bfseries\addfontfeature{LetterSpace=6}BUDDIMM · STAGE IA AGENTIQUE 2027 · NOTE DE RENDU DU TEST SITADEL}\\[3pt]
{\color{encre}\fontsize{21}{24}\selectfont\bfseries Le demandeur n'est pas le client}\\[2pt]
{\color{gris}\normalsize Qualifier qui construit dans les Yvelines, retrouver qui exploitera le bâtiment, et en tirer une liste de prospects sourcée}\par}
\vspace{6pt}

\noindent\begin{tikzpicture}
  \node[fill=fond,draw=none,inner xsep=9pt,inner ysep=7pt,text width=\dimexpr\linewidth-18pt\relax,
        path picture={\fill[accent] (path picture bounding box.north west) rectangle ([xshift=2.2pt]path picture bounding box.south west);}] {%
  \sffamily\raggedright
  \begin{minipage}[t]{42mm}\champ{CANDIDAT}{\Large\bfseries\color{encre}BAQUA Abdellah}\end{minipage}\hfill
  \begin{minipage}[t]{24mm}\champ{TEMPS PASSÉ}{\Large\bfseries\color{encre}7 heures}\end{minipage}\hfill
  \begin{minipage}[t]{57mm}\champ{DÉPÔT GIT}{\small\bfseries\color{accent}\href{https://github.com/AbdellahBaqua/SitadelAgent}{github.com/AbdellahBaqua/SitadelAgent}}\end{minipage}\hfill
  \begin{minipage}[t]{28mm}\champ{RENDU LE}{\small\bfseries\color{encre}25 novembre 2026}\end{minipage}\par
  \vspace{6pt}
  \champ{OUTILS UTILISÉS}{\footnotesize\color{encre}Claude Code (assistant de code) · NVIDIA NIM (Nemotron 3 Ultra et Lightning) · Groq (gpt-oss-120b) · Python · \LaTeX}\par
  \vspace{4pt}
  {\scriptsize\color{gris}Données publiques uniquement (SDES, Annuaire des entreprises, FINESS) · personnes morales seulement, aucun nom de personne physique · aucune prise de contact avec les demandeurs\par}};
\end{tikzpicture}

\vspace{7pt}
\noindent\kpi{encre}{3\,957}{demandeurs personnes morales\\(3\,363 PC logements, 4\,245 PC locaux)}\hfill
\kpi{accent}{92\,\%}{classés par des règles du code ;\\8\,\% par l'agent, 58 non déterminés}\hfill
\kpi{orange}{22}{exploitants retrouvés dans FINESS\\sur 98 permis, dont 8 distincts\\du demandeur}\hfill
\kpi{vert}{131}{prospects jugés par l'agent,\\après 99 écartés par le code}

% ================= 1. angle =================
\section{L'angle choisi, et pourquoi}
Sitadel dit qui \emph{dépose} un permis, pas ce qu'est le demandeur ni qui \emph{exploitera} le bâtiment. Or Buddimm accompagne
des associations médico-sociales, des structures de l'ESS et des bailleurs sociaux : la bonne question n'est pas seulement
\og quelle famille ?\fg{}, mais \og derrière ce permis, quelle structure a un projet immobilier et pas de direction immobilière ?\fg{}
Le prototype y répond avec un principe constant : \textbf{le code tranche quand un champ fiable suffit, l'agent LLM ne traite que
ce qui demande un jugement, et le code contrôle l'agent} (liste fermée, source citée pour chaque indice, fiche entreprise
obligatoire, plafonds de confiance, trace complète). Périmètre : PC de logements et de locaux des Yvelines (millésime 2026-10)
dont le demandeur est une personne morale.

\vspace{3pt}
\begin{center}
\begin{tikzpicture}[node distance=8mm and 7.5mm]
  \node[etape=gris] (sit) {\textbf{Sitadel 78}\\PC logements + locaux\\personnes morales};
  \node[etape=accent,right=of sit] (reg) {\textbf{Règles (code)}\\catégorie juridique,\\NAF, dénomination};
  \node[etape=orange,right=of reg] (ag1) {\textbf{Agent classification}\\fiche Annuaire,\\libellés Insee};
  \node[etape=vert,right=of ag1] (fam) {\textbf{8 familles}\\+ \og non déterminé\fg{}\\lien source par ligne};
  \draw[fleche] (sit) -- (reg);
  \draw[fleche] (reg) -- node[above=1pt,chiffre=accent]{8\,\%} (ag1);
  \draw[fleche] (ag1) -- (fam);
  \node[etape=gris,below=of sit] (tie) {\textbf{Projets médico-sociaux}\\déposés par un tiers\\(promoteur, bailleur)};
  \node[etape=orange,right=of tie] (ag2) {\textbf{Agent exploitant}\\candidats FINESS\\dans la commune\\et la rue};
  \node[etape=accent,right=of ag2] (pre) {\textbf{Pré-filtre (code)}\\culte, école, sport,\\cabinet, structure cessée};
  \node[etape=orange,right=of pre] (ag3) {\textbf{Agent qualification}\\Annuaire + FINESS\\+ garde-fou preuve dure};
  \draw[fleche] (tie) -- (ag2);
  \draw[fleche] (ag2) -- node[above=1pt,chiffre=orange]{+6} (pre);
  \draw[fleche] (pre) -- node[above=1pt,chiffre=accent]{$-$99} (ag3);
  \draw[fleche] (fam.south) -- node[right=1pt,font=\sffamily\scriptsize,color=gris]{230 candidats} ($(fam.south)+(0,-3.8mm)$) -| (pre.north);
  \node[etape=vert,below=4.5mm of ag3] (pro) {\textbf{Prospects priorisés}\\score explicable, preuves\\et sources par ligne};
  \draw[fleche] (ag3) -- (pro);
  \begin{scope}[on background layer]
    \node[fit=(sit)(fam)(pro),draw=filet,dashed,rounded corners=3pt,inner sep=3.5pt] (cadre) {};
  \end{scope}
  \node[anchor=south west,font=\sffamily\scriptsize,color=gris,text width=98mm,align=left] at ($(cadre.south west)+(4pt,3pt)$)
    {\textcolor{accent}{\rule{6pt}{6pt}} code déterministe \quad \textcolor{orange}{\rule{6pt}{6pt}} agent LLM encadré \quad \textcolor{vert}{\rule{6pt}{6pt}} sortie sourcée\\[1pt]
     garde-fous : liste fermée · source citée · fiche obligatoire · plafonds · traces};
\end{tikzpicture}
\end{center}
\vspace{-6pt}

% ================= 2. résultats =================
\section{Ce que j'ai trouvé}
\begin{itemize}
\item \textbf{Les associations construisent rarement.} Côté logements, les 30 demandeurs associatifs ou exploitants santé-social ne déposent que
  47 PC ; seuls 6 en déposent plus d'un. Un permis est donc un signal \emph{tardif et ponctuel} : le compte se gagne avant le dépôt.
\item \textbf{Le médico-social est déposé par d'autres.} Sur 64 PC de logements en résidence médico-sociale ou sociale, 49 (77\,\%) viennent de
  promoteurs (26) et de bailleurs (23), 4 seulement d'associations. D'où l'\emph{agent exploitant}, qui rapproche ces permis de
  FINESS (registre des établissements sanitaires et sociaux) : sur 98 permis traités, 22 exploitants trouvés, dont 8 distincts du
  demandeur, par exemple l'association \emph{Handi Val de Seine} derrière une résidence déposée par le bailleur Logement Francilien.
\item \textbf{Le vrai promoteur est masqué.} 42\,\% des PC de logements des promoteurs passent par des sociétés de construction-vente à usage unique ;
  en lisant leurs dirigeants personnes morales, 633 des 795 sociétés de programme sont rattachées à un groupe.
\item \textbf{Les équipements de santé portés par des associations sont rares mais lourds} : 33,5\,\% de la surface de santé créée
  (79\,217\,m²), dont un seul projet pèse 46\,\%. Une poignée de comptes concentre l'enjeu.
\item \textbf{Une liste de prospects nettoyée.} Sur 230 candidats des familles cibles, le code écarte 99 cas sans ambiguïté
  (culte, écoles, clubs sportifs, cabinets libéraux, structures cessées) ; l'agent ne juge que les 131 restants. Tête de liste :
  Avenir APEI, Fondation Mallet, Delos APEI 78, Autisme en Yvelines, Fondation John Bost, Handi Val de Seine, Coallia, Vilogia.
\end{itemize}

% ================= 3. preuves =================
\clearpage
\section{Exemples : prospects qualifiés et leurs preuves}
{\small
\begin{tabularx}{\linewidth}{@{}P{29mm}P{54mm}Y@{}}
\toprule
\textsf{\footnotesize\bfseries\color{gris}PROSPECT} & \textsf{\footnotesize\bfseries\color{gris}PREUVE DURE (VÉRIFIÉE PAR CODE)} & \textsf{\footnotesize\bfseries\color{gris}BESOIN IMMOBILIER VISIBLE (AGENT, SOURCES CITÉES)}\\
\midrule
\textbf{Avenir APEI} & 20 établissements FINESS ; activité 88.10B (accompagnement d'adultes handicapés) & 3 permis autorisés en 2024--2025, pas encore commencés : 8 logements pour personnes handicapées au Vésinet, 564 et 576\,m² d'action sociale à La Celle-Saint-Cloud.\\
\textbf{Fondation Mallet} & 9 établissements FINESS ; activité 86.10Z (hospitalier) & 34 permis de 2014 à 2025, chaque année, autour de Houdan, Richebourg et Bazainville : un maître d'ouvrage régulier, rare pour une fondation.\\
\textbf{Handi Val de Seine} & 18 établissements FINESS ; activité 87.30B ; reliée par l'agent exploitant & Exploite la résidence de 30 logements de Follainville-Dennemont, déposée par le bailleur Logement Francilien : seul FINESS relie ce bâtiment à son exploitant.\\
\textbf{Coallia Habitat} & Bailleur social ; permis de résidence sociale & Résidence sociale de 115 logements autorisée en 2025 à Sartrouville, pas encore commencée.\\
\bottomrule
\end{tabularx}}
\vspace{1pt}{\footnotesize\color{gris}Chaque ligne de \texttt{prospects\_priorises.csv} porte ces colonnes (\texttt{preuves\_dures}, \texttt{indices\_agent}, \texttt{raisons} du score) et le lien vers sa fiche publique ; les 99 candidats écartés gardent leur règle dans \texttt{prospects\_ecartes.csv}.\par}

% ================= 4. établi / fragile / manque =================
\section{Ce qui est établi, fragile, ou manque}
\noindent\bilan{vert}{Établi}{Familles par catégorie juridique Sirène (92\,\% des demandeurs, exact par construction). Tous les comptages,
  rejouables depuis les données. Chaque ligne renvoie à sa source publique (SDES, Annuaire des entreprises, FINESS).
  Les 99 exclusions du pré-filtre, relues une à une.}\hfill
\bilan{orange}{Fragile}{Le jugement de l'agent sur les cas ambigus : deux modèles s'accordent sur 56\,\% des classifications difficiles,
  et sur 20 des 22 qualifications de prospects \emph{après} le garde-fou \og preuve dure\fg{} (17/22 avant). Le rapprochement FINESS :
  40\,\% des cas où l'exploitant devrait être connu. Le barème du score et les seuils de taille.}\hfill
\bilan{rouge}{Manque}{Une précision mesurée : les échantillons de contrôle sont tirés (50 demandeurs, 12 verdicts) mais pas encore
  vérifiés à la main. La qualification complète (22 des 131 prospects, en test). Les successeurs des bailleurs fusionnés.
  Les projets pas encore ouverts, invisibles dans FINESS.}

% ================= 5. ce qui bloque =================
\section{Ce qui bloque}
\begin{itemize}
\item \textbf{Le NAF d'un permis est daté} : c'est celui d'un établissement au dépôt, souvent fermé depuis. L'agent s'appuie sur
  l'activité \emph{actuelle} de l'Annuaire et signale toute divergence.
\item \textbf{Les champs utiles sont mal remplis} : le type de résidence est \og non rempli\fg{} (code 9) pour 41 des 47 permis de
  logements associatifs (87\,\%) ; \og résidence pour personnes âgées\fg{} mêle Ehpad et résidences services seniors privées, absentes
  de FINESS ; environ 15\,\% des ouvertures de chantier et un tiers des achèvements ne remontent jamais.
\item \textbf{La cible n'est pas définie assez finement.} Les désaccords entre modèles portent presque tous sur des associations
  \og grises\fg{} (culturelles, communautaires, de formation) : est-ce de l'ESS au sens de Buddimm ? C'est une question métier, pas
  de puissance de modèle ; le garde-fou du code la tranche provisoirement en \og à vérifier\fg{}.
\item \textbf{Une erreur que le code ne voit pas} : les constats rédigés par les agents explorateurs ont des chiffres vérifiés
  automatiquement, mais 3 interprétations sur 19 étaient fausses ; je les ai retirées après relecture.
\end{itemize}

% ================= 6. ensuite =================
\section{Ce que je ferais ensuite}
\begin{enumerate}[label=\textcolor{accent}{\sffamily\bfseries\arabic*},leftmargin=1.2em,itemsep=2.2pt]
\item Vérifier à la main les échantillons de contrôle et publier la précision des règles et de l'agent, strate par strate.
\item Fixer avec l'équipe la définition de la cible (crèches, associations culturelles, très grands groupes) et en faire des règles.
\item Qualifier les 131 prospects avec le modèle Groq : même qualité que le modèle NVIDIA sur l'échantillon, 4 fois plus rapide (17\,s contre 73\,s par cas).
\item Améliorer le rapprochement FINESS (adresses abrégées, communes seules) et retrouver le successeur des bailleurs fusionnés.
\item \textbf{Brancher le résultat sur la prospection}, chaque mois :
\end{enumerate}
\vspace{3pt}
\begin{center}
\begin{tikzpicture}[node distance=3mm and 5mm]
  \tikzset{petit/.style={etape=#1,minimum height=9mm,text width=26mm}}
  \node[petit=gris] (a) {\textbf{Sitadel}\\mise à jour mensuelle};
  \node[petit=accent,right=of a] (b) {\textbf{Nouveaux permis}\\et changements d'état};
  \node[petit=orange,right=of b] (c) {\textbf{Pipeline}\\famille, exploitant,\\qualification};
  \node[petit=vert,right=of c] (d) {\textbf{Fiche prospect}\\besoin, permis, sources};
  \node[petit=accent,right=of d] (e) {\textbf{Revue humaine}\\puis CRM};
  \foreach \x/\y in {a/b,b/c,c/d,d/e} \draw[fleche] (\x) -- (\y);
\end{tikzpicture}
\end{center}
\vspace{-5pt}
Les permis nouveaux ou passés à l'état \og autorisé, non commencé\fg{} déclenchent une fiche : qui est la structure, quel projet,
quel besoin probable, avec les liens vers les sources publiques. Alertes utiles : un exploitant connu dépose à nouveau, un bailleur
dépose une résidence médico-sociale sans exploitant identifié, un permis associatif est annulé (projet à sécuriser).
\textbf{La décision de contacter reste humaine, et le contact n'est jamais automatisé.}

\end{document}
